% Body-only supplement.  Requires amsmath, amssymb, booktabs and hyperref.
% All counts are generated by the adjacent standard-library Python programs.
\section{Measured encoding changes and final finite-theory controls}
\label{sec:encoding-controls}

The six-function FC13S and FC14S theories admit two further, rigorous
simplifications: one class-level Boolean cancellation and one explicitly
charged selector constant.  These produce smaller finite presentations,
\emph{not a compiled binary Turing machine}.  The most straightforward
two-constant target is FC14Qd, a definitional extension of the Boolean-normalized
raw closure theory.  Its native stream has 926 bits and its equality-completed
stream has 1401 bits in the same four-bit term-node format.  The corresponding
extensional control FC15Qd has 1085 and 1560 bits.  State counts, decoder costs,
blank-tape initialization, and a machine-level proof remain unmeasured.

\subsection{A class-level Boolean cancellation: FC13B and FC14B}

Write $c$ for the previously specified closedness class and $s$ for its support
class.  The S presentation uses the hull class
\[
 \operatorname{Ex}(r,\operatorname{And}(c,s))
 =I\bigl(U(r,I(I(I(c,I(s,d)),d),d)),d\bigr).
\]
The B presentation replaces this class expression by
\[
 \boxed{I\bigl(U(r,I(c,I(s,d))),d\bigr).}
\]
No other clause changes.  In particular, the freezing equation, the witness
coordinate, the support condition, and the selector guards are unchanged.

\paragraph{Core-relative equivalence.}
The implication and falsity clauses give, for every assignment $a$ and class
object $x$,
\[
 \operatorname{Inc}(a,I(I(x,d),d))
 \ \longleftrightarrow\ \operatorname{Inc}(a,x).
\]
The exact universal-class semantics, obtained from the universal clause and
its counterexample clause, preserves pointwise incidence equivalence in its
class argument.  So does $I$ in each of its class arguments.  Apply these facts
to $x=I(c,I(s,d))$ inside the displayed $U$ argument and then its surrounding
$I$ expressions.  The old and new closure clauses are equivalent in every
model of the common core.  Thus FC13B has exactly the same models as FC13S in
the same signature; the same statement holds for FC14B and FC14S.  This is
\emph{not} an equality of the underlying class objects.  The replacement is
not performed inside an assignment coordinate, a selector label, or an
arbitrary argument of another function.

The cancellation removes two $I$ nodes and two $d$ leaves: exactly four term
nodes and sixteen flat bits, native and completed.  The maximum template depth
drops from 15 to 13.  The executable identity control checks all 680 choices
of antecedent truth and two class truth tables on fibers of sizes one through
four.  These are Boolean identity checks, not finite models of the full core.

\subsection{Name the selector: FC14Q/Qd and FC15Q/Qd}

Starting with a B theory, replace each designated selector term $I(d,d)$ by a
new constant $q$.  The proper function signature remains
$V,C,I,A,U,F$.  There are two alternatives:
\[
 \begin{array}{ll}
 \mathrm{Q}:& \operatorname{Inc}(a,q)\quad\hbox{for all }a,\\
 \mathrm{Qd}:&q=I(d,d).
 \end{array}
\]
The Q axiom costs fourteen flat bits; Qd costs twenty-two.  Both add one clause
and one literal.  There is no extra positive-arity congruence clause for a new
constant: universal equality reflexivity already supplies $q=q$, and a
constant denotes a single element in the quotient.

\paragraph{The definitional version.}
FC14Qd and FC15Qd are literal definitional extensions of FC13B and FC14B.
Expanding the new constant recovers the B clauses; conversely any B model
expands by interpreting $q$ as $I(d,d)$.  This is the cleanest additional
consistency endpoint.  Equality completion is applied afterward, using the
same retained-language quotient argument as before.

\paragraph{The truth-only version.}
The Q theories are also equiconsistent with their B counterparts by explicit
native-equality model transports.  A B model expands with $q=I(d,d)$, which is
incident to every assignment because $d$ is incident to none.  Conversely,
in a Q model change only the implication function at the input pair $(d,d)$:
\[
 I'(x,y)=
 \begin{cases}q,&x=d\text{ and }y=d,\\ I(x,y),&\text{otherwise.}\end{cases}
\]
Its incidence semantics is still exact implication.  In the Q presentation,
all occurrences of the designated selector are the constant $q$, so the
assignment, atomic-class, freezing, and counterexample-witness equations are
unchanged.  A structural induction through the closure and extensionality
class terms, confined to their $I/U$ class positions, shows that these clauses
retain their incidence truth.  Now $I'(d,d)=q$, and forgetting the name $q$
gives a B model.  This argument does not perform a piecewise change on an
arbitrary prequotient equality predicate; genuine equality is used first, and
completed presentations are related by the existing quotient construction.

There are fifteen designated-selector occurrences in raw B and nineteen in
extensional B.  Replacing each three-node term by a one-node constant saves
eight bits per occurrence before the added axiom is charged.

\begin{center}
\begin{tabular}{lrrrrrr}
\toprule
Target & Constants & Native C/L & Complete C/L & Native bits & Complete bits & Depth\\
\midrule
FC13S &1&13/26&23/57&1040&1515&15\\
FC13B &1&13/26&23/57&1024&1499&13\\
FC14Q &2&14/27&24/58&918&1393&12\\
FC14Qd&2&14/27&24/58&926&1401&12\\
\midrule
FC14S &1&14/30&24/61&1231&1706&15\\
FC14B &1&14/30&24/61&1215&1690&13\\
FC15Q &2&15/31&25/62&1077&1552&12\\
FC15Qd&2&15/31&25/62&1085&1560&12\\
\bottomrule
\end{tabular}
\end{center}
Here C/L means clauses/literals.  Every row has six proper functions and two
non-equality predicates.  Native clauses use at most four local variables;
completion uses six.  Flat bits include three-bit clause lengths, three-bit
literal headers, every expanded four-bit term node, and the final three-bit
terminator.  A fixed known signature is assumed, exactly as in the source44
comparison; these data counts do not include a signature decoder program.

\subsection{Linear-size ground-term codes}

The original radix-eight Cantor reference is convenient for a total natural
number decoder, but its bit lengths can grow sharply with nested terms.
The alternative ground codec assigns fixed three-bit prefix tags to the
retained function and constant symbols.  With one constant the ordered
alphabet is $d,V,C,I,A,U,F$; with two it is $d,q,V,C,I,A,U,F$.
The known arities are $0$ for constants, $3$ for $C$, and $2$ for the other
proper functions.  No variable opcodes occur in these \emph{ground} words.

\paragraph{Unique reconstruction and exact cost.}
Start with one pending child slot.  Reading a symbol of arity $k$ consumes one
slot and creates $k$.  A valid word reaches zero exactly at its final token
and never earlier.  Parsing in prefix order therefore reconstructs a unique
ordered ground tree, and induction on that tree proves that parsing and
encoding are inverse.  Its length is exactly $3\lvert t\rvert$ bits, where
$\lvert t\rvert$ counts all constructor and constant nodes.  If a natural
number name is required, prepend one leading $1$: its binary bit length is
$3\lvert t\rvert+1$, and leading-zero tokens are preserved.  The implementation
uses strings as atom-name components, avoiding arithmetic pairing of those
already large numbers.

\paragraph{Fair enumeration and the halting endpoint.}
Enumerate finite token words by length and then lexicographic tag order,
discarding words failing the shape check.  Each finite block is finite, every
retained ground tree has one valid word, and each occurs exactly once.  The
prefix of the resulting term list is a monotone substitution frontier.
At stage $n$, substitute every tuple from that frontier into every completed
clause, allowing repeated values and retaining all endpoint terms outside the
frontier.  Decide that finite propositional ground CNF exactly.  Halt only
after an exhausted UNSAT decision; after SAT, enlarge the frontier.

Every finite set of retained ground instances occurs by some finite stage.
Compactness and the equality quotient therefore give the same ideal
inconsistency semidecision as the original reference.  Injective renaming of
ground terms preserves the finite propositional clause set up to atom names.
Combining this fact with the finite-theory consistency calibrations gives, in
ordinary ambient ZFC, the requested mathematical endpoint for the ideal
refuter.  This is not a serialized PA proof and not a machine compilation.
Malformed token words are skipped; an exception, memory exhaustion, or
recursion limit is never converted into UNSAT.  The reference implementation
has no arbitrary word-length cutoff.

For a fair concrete comparison, all following rows use the \emph{same}
substitution frontier $\{d,V(d,d)\}$ and completed clauses:
\begin{center}
\begin{tabular}{lrrrrr}
\toprule
Target & Ground C/A & Max. nodes & Prefix bits & Integer bits & Cantor bits\\
\midrule
FC13S &300/259&115&345&346&63197\\
FC14S &308/267&115&345&346&63197\\
FC13B &300/259&111&333&334&15798\\
FC14B &308/267&111&333&334&15798\\
FC14Q &302/261&95&285&286&11796\\
FC14Qd&301/260&95&285&286&11796\\
FC15Q &310/269&95&285&286&11796\\
FC15Qd&309/268&95&285&286&11796\\
\bottomrule
\end{tabular}
\end{center}
Ground C/A means clauses/atoms.  These are separate finite SAT instances,
each with an independently reevaluated SAT witness.  Within each row the
prefix and numeric representations produce exactly the same numbered CNF.
Across different theories no assertion of identical CNFs is made.  For the
two-constant comparator, numeric names are $d=0$, $q=1$, and
$2+i+8\operatorname{tuple}(\vec x)$ for function tag $i$; unused radix tags
are unary padding functions used only by that arithmetic comparator.

The two-constant prefix enumerator's \emph{own} first two terms are $d,q$,
not $d,V(d,d)$.  On that different frontier the Q/Qd maximum endpoint has
71 nodes, or 213 prefix bits.  Those smaller figures must not be substituted
into the common-frontier table.  Both alphabets have been checked against an
independent slot-count validator for all 37448 three-bit token words of one
through five tokens; the two-constant numeric codec also has 4096 roundtrips.

This change removes the exhibited representational blowup in ground-name
length.  It does not remove the number of substitutions, SAT search, or
potentially slow valid-word enumeration.  Prefix decoding, stack handling,
term copying, and the initial finite theory still need implementations in any
binary machine.  No numerical bound on those control states follows from
the displayed word lengths.

\subsection{Fully charged shared-template experiments}

The separate source-template experiment reconstructs syntax, including local
variable leaves.  It does not add semantic class constructors or clauses.
The original source44 FC15*/FC16* trees are stored in a portable snapshot with
the ProveIt commit, generator paths, and SHA-256 hashes.  The same codec rules
and searches apply to old and new theories.

Three variable naming policies are tested: lexical source names; first
occurrence in the fixed clause/literal/tree order; and one forward greedy
pass through clauses.  In the last policy all $n!$ bijections from a clause's
$n$ variables onto $v_0,\ldots,v_{n-1}$ are examined, choosing the one that
adds the fewest distinct constructor subtrees to the already chosen clauses;
lexical permutation order breaks ties.  There is no global-optimality claim.
No clauses or literals are reordered.  The reconstructed payload is exactly
the chosen alpha-renamed syntax, and stored audit maps invert it to the
original syntax.  Those maps are not required by the universal theory and
are not counted as payload.

\paragraph{Full constructor DAG.}
The fixed terminal dictionary contains six variable slots followed by the
declared constants.  Every distinct proper constructor subtree is emitted,
ordered by increasing depth and then a deterministic textual tie-breaker.
If there are $m$ constructor records, the stream starts with Elias gamma
encoding of $m+1$.  A record gives the constructor opcode and references to
earlier terminals or records; if $b$ values are available, each reference has
$\lceil\log_2 b\rceil$ bits.  The clause body retains the original clause
and literal headers, with two final-dictionary references per literal, plus
the terminating clause-size zero.  Every reference, header, and record is
charged.  Decoding forbids forward references and reconstructs each tree.

\paragraph{Sparse subtree grammar.}
Only selected repeated subtrees are named.  The dictionary count again uses
gamma$(m+1)$ and entries have arity-delimited tree encodings, referring only
to earlier macros.  Let $w$ be the ordinary term-opcode width and reserve the
all-ones $w$-bit opcode as an escape.  Escape followed by $0$ denotes that
original opcode, if present; escape followed by $1$ denotes a macro reference
with $\lceil\log_2 k\rceil$ index bits when $k$ macros are available.
Other opcodes retain their original $w$ bits.  Thus in the original source44
signature its highest constructor $H$ pays a five-bit opcode instead of four;
the exact cost is included.  A macro reference is used only if strictly
shorter than its expansion at that point.  The fixed known signature, rather
than an uncharged per-theory compression dictionary, determines all arities.

Candidate macros are precisely proper constructor subtrees occurring at least
twice in the chosen syntax.  Start with no macros; repeatedly examine every
single add/remove move, recomputing the \emph{whole} charged bitstream, and
accept the greatest strict improvement.  Fixed depth/text order breaks ties.
The integer bit cost decreases at every accepted step.  This terminates but
does not prove a global minimum or account for decoder state cost.

The table gives the best value from the three declared variable naming
policies for each codec.  All input signatures here use the same fixed six
variable slots and four-bit term opcodes.
\begin{center}
\begin{tabular}{lrrrrrr}
\toprule
&\multicolumn{3}{c}{Native}&\multicolumn{3}{c}{Equality completed}\\
Target & Flat & Full DAG & Sparse & Flat & Full DAG & Sparse\\
\midrule
FC15* &881&900&828&1616&1870&1543\\
FC13S &1040&1244&927&1515&1938&1383\\
FC13B &1024&1212&911&1499&1792&1367\\
FC14Q &918&1220&866&1393&1800&1322\\
FC14Qd&926&1237&874&1401&1816&1330\\
\midrule
FC16* &1088&1190&1020&1823&2090&1729\\
FC14S &1231&1519&1094&1706&2173&1569\\
FC14B &1215&1483&1078&1690&2137&1553\\
FC15Q &1077&1494&1014&1552&2149&1475\\
FC15Qd&1085&1512&1022&1560&2168&1483\\
\bottomrule
\end{tabular}
\end{center}
Full DAG naming loses against flat syntax throughout: repeated small nodes
do not repay dictionary and root-reference costs.  Sparse naming does help.
For the native S/B targets its selected macros are just the fixed truth
selector and the repeated assignment pullback in the two freezing clauses.
The deep closure expression is mostly unrepeated; giving it a macro name
alone does not make it cheap.  A reduced proper-function signature chiefly
helps the equality-completed presentation in this experiment.

\paragraph{Apply the constant idea to prior art too.}
For a fair control, name source44's $N(d,d)$ by a constant $q$ and add
$q=N(d,d)$.  One must price the changed template alphabet.  Six variable
slots plus two constants plus nine functions give seventeen opcodes: the
fixed-slot flat format now needs five bits.  This threshold is not an
unavoidable native lower bound.  Native source44 clauses use only four
variables; reclaiming the unused slots restores a four-bit alphabet.
\begin{center}
\begin{tabular}{lrrrr}
\toprule
Prior control & Variable slots & Opcode bits & Flat bits & Best sparse bits\\
\midrule
FC15* Qd native&6&5&1029&963\\
FC15* Qd native&4&4&855&811\\
FC15* Qd complete&6&5&1902&1805\\
FC16* Qd native&6&5&1234&1161\\
FC16* Qd native&4&4&1022&975\\
FC16* Qd complete&6&5&2107&1986\\
\bottomrule
\end{tabular}
\end{center}
Consequently the new Qd targets do not beat the optimized old native payloads.
Their useful measured advantage is in the completed presentations: compare
1401/1560 flat or 1330/1483 sparse bits with the old best original completed
1616/1823 flat or 1543/1729 sparse streams.  A future compiler that handles
native equality directly may rank these alternatives differently.  Taking
different best encodings per theory is a cost comparison, not evidence that
one common machine can decode all of them without additional control.

\subsection{Rejected three-function channel encoding}

One further control replaces $I,A,U$ by values stored at coordinate $d$ of
already required $F$ outputs, leaving $V,C,F$.  With the truth-only named
selector, let $P(a,b)=C(C(d,d,a),q,b)$ and choose pairwise distinct tags
\[
 \tau_I=P(d,d),\quad \tau_A=P(d,q),\quad \tau_U=P(q,d).
\]
Replace a constructor $G(a,b)$ for $G\in\{I,A,U\}$ by
\[
 V\bigl(F(P(\tau_G,P(a,b)),d),d\bigr).
\]
The lower term interpretation is immediate if the resulting native theory
has a model.  A proposed upper construction would extend rich class coding
to every two-coordinate fiber: coordinate $d$ stores the desired constructor
value, while coordinate $q$ retains the quantifier witness.  Distinct tagged
keys prevent conflicting value requests, and a third coordinate would make
each such fiber infinite.  This is only a proof sketch in this supplement;
the control is not added to the fully proved recommended endpoints.

Its exact syntactic price already rejects it for the tested formats:
\begin{center}
\begin{tabular}{lrrrr}
\toprule
Control & Clauses/literals & Term nodes & Flat bits & Lexical sparse bits\\
\midrule
Raw native &14/27&978&4038&2625\\
Raw complete&21/49&1036&4357&2944\\
Extensional native&15/31&1234&5077&3306\\
Extensional complete&22/53&1292&5396&3625\\
\bottomrule
\end{tabular}
\end{center}
Every row has depth 56.  Removing three congruence clauses does not compensate
for the expanded coordinate-address terms.  No state bound or historical
novelty claim is inferred from this negative experiment.

\subsection{Portable reproduction and scope}

The new programs use only Python's standard library and adjacent source files.
Run with assertions enabled:
\begin{verbatim}
python3 boolean_class_control.py
python3 linear_ground_codec.py
python3 linear_ground_codec.py --boolean-normalized \
  --output linear-ground-boolean-results.json
python3 shared_template_codec.py
python3 named_selector.py
python3 channel_collapse_control.py
\end{verbatim}
The original \texttt{frozen\_kernel.py} and \texttt{frozen\_frontend.py}
reproduce the unchanged S controls.  Each new payload has an independent
decoder, exact tree reconstruction checks, and truncation/trailing-data
negative controls.  The emitted \texttt{.shared.json} files record the
variable maps and selected macro trees.  The \texttt{.bits} files contain
the precise streams counted in the tables.

Optional ideal unbounded searches are
\begin{verbatim}
python3 linear_ground_codec.py --whole-search
python3 linear_ground_codec.py --boolean-normalized --whole-search
python3 named_selector.py --whole-search
python3 named_selector.py --truth-only --whole-search
\end{verbatim}
Add \texttt{--extensional} for the extensional target.  These commands are
executable reference refuters, with no expectation of practical termination.
They should not be used as quick reproduction checks.  The searches halt
successfully only after a finite UNSAT result; ordinary host failures
propagate.  Their consistency endpoint is the ordinary-ZFC endpoint proved
for the finite theories, rather than a new verified machine-level result.

The resulting compiler choices are concrete: FC13B keeps one constant and
1024 native bits; FC14Qd pays one explicit definition clause to reduce the
raw completed flat stream to 1401 bits; FC15Qd supplies the corresponding
extensional control.  Truth-only Q saves another eight bits at the cost of
the additional native model-transport argument.  Prefix ground names avoid
the exhibited Cantor bit-length growth, while sparse template sharing offers
separate data savings.  A binary compiler must now measure those choices
together with its decoder and proof-search control.
